\section{ERR-010: JSON Parse Errors in Tool Output}
\label{sec:err-010}
\subsection*{Metadata}
\begin{tabular}{ll}
\textbf{Issue ID} & ERR-010 \\
\textbf{Severity} & MEDIUM \\
\textbf{Category} & Error Handling \\
\textbf{Branch} & \branchref{fix/ERR-010-json-parse-errors}
\textbf{Changes} & \branchdiff{fix/ERR-010-json-parse-errors} \\
\end{tabular}

\subsection*{Executive Summary}
In \srcfile{src/tool-loader.ts}, after a declarative tool script exits successfully, the code attempts to parse its stdout as JSON. If the output starts with \code{\{} but is not valid JSON — truncated output, log prefixes, or trailing content — \code{JSON.parse()} will throw and the entire output is treated as raw text, potentially losing structured data. More critically, if stdout is empty or whitespace-only, the empty string is replaced with ``(no output)'' which is misleading.

\subsection*{Root Cause Analysis}
The original code attempts JSON parsing unconditionally:

\begin{lstlisting}[language=TypeScript]
// BEFORE — no guard before JSON.parse
let text = stdout.trim();
try {
    const parsed = JSON.parse(text);
    if (parsed.text) text = parsed.text;
    else if (parsed.result) text = typeof parsed.result === "string" ? parsed.result : JSON.stringify(parsed.result);
} catch {
    // Raw text output is fine
}

resolve({
    content: [{ type: "text", text: text || "(no output)" }],
    details: undefined,
});
\end{lstlisting}

If stdout is partial JSON with trailing content (e.g., \code{\{"text": "hello"\}\textbackslash nsome-debug-log}), \code{JSON.parse} throws and the entire string becomes raw text.

\subsection*{Fix Implementation}
A guard checks for JSON-like output before parsing:

\begin{lstlisting}[language=TypeScript]
// AFTER — guard before JSON.parse
let text = stdout.trim();
if (text && (text.startsWith("{") || text.startsWith("["))) {
    try {
        const parsed = JSON.parse(text);
        if (parsed.text) text = parsed.text;
        else if (parsed.result) text = typeof parsed.result === "string" ? parsed.result : JSON.stringify(parsed.result);
    } catch {
        // Raw text output is fine
    }
}

resolve({
    content: [{ type: "text", text: text || "(no output)" }],
    details: undefined,
});
\end{lstlisting}

\subsection*{Affected Files}
\begin{itemize}
    \item \srcfile{src/tool-loader.ts} — Lines 119-132
    \item \srcfile{src/\_\_tests\_\_/json-parse-errors.test.ts}
\end{itemize}

\subsection*{Verification}
Tests verify that non-JSON output is not parsed:

\begin{lstlisting}[language=TypeScript]
const text = "some plain text output";
const shouldParse = text && (text.startsWith("{") || text.startsWith("["));
assert.equal(shouldParse, false);

const jsonText = '{"text": "hello world"}';
const shouldParseJson = jsonText && (jsonText.startsWith("{") || jsonText.startsWith("["));
assert.equal(shouldParseJson, true);
\end{lstlisting}
